% Part: history
% Chapter: biographies 
% Section: kurt-goedel

\documentclass[../../../include/open-logic-section]{subfiles}

\begin{document}

\olfileid{his}{bio}{god}

\olsection{کرټ ګوډل}

\olphoto{goedel-kurt}{کرټ ګوډل}

کرټ ګوډل (\textsc{ګېړ}-ډل) د 1906 د اپرېل پر
28مه د اتريش--هنګري په امپراتورۍ کښې په برون
کښې زېږېدلے و (چې اوس د چېک جمهوريت برنو دے)۔
د ځيرک او پوښتونکي طبيعت له امله به ئې کورنۍ
ځوان کرټلې ته ډېر ځله ``Der kleine Herr Warum''
(کوچنے ښاغلے ولې) ويل۔ هغه له لومړني ښوونځي
څخه په زده کړو کښې غوره و، او يوازې په
رياضياتو کښې ئې کله تر ټولو لوړې درجې کمه
درجه اخيسته۔ ګوډل د خرابې روغتيا له امله
ډېر ځله له ښوونځي غيرحاضر و او له بدني روزنې
څخه مستثنا و۔ په ماشومتوب کښې ئې د روماتيکي
تبې تشخيص شوے و۔ د ټول ژوند په اوږدو کښې ئې
دا باور درلود چې دې تبې ئې زړه د تل دپاره
اغېزمن کړے دے، که څه هم طبي ارزونې د دې
خلاف ويل۔

ګوډل په 1924 کښې د ويانا په پوهنتون کښې زده
کړه پيل کړه او په~1929 کښې ئې د دوکتورا زده
کړې بشپړې کړې۔ لومړے ئې د فزيک د زده کړې
اراده وه، خو شوق ئې ژر رياضياتو او په ځانګړي
ډول منطق ته واوښت، چې د فيلسوف روډولف کارناپ
اغېز هم پکې برخه لرله۔ د هانس هان تر
لارښوونې لاندې ليکلې رسالې ئې د عينيت لرونکي
لومړي ترتيب د پريديکات منطق د بشپړتيا قضيه
ثابته کړه \citep{Godel1929}۔ يوازې يو کال
وروسته ئې خپلې تر ټولو مشهورې نتيجې ترلاسه
کړې: د ناتکميلۍ لومړۍ او دويمه قضيه
(په \citealt{Godel1931} کښې خپرې شوې)۔ په
ويانا کښې د اوسېدو پر وخت ګوډل د ويانا له
کړۍ سره ډېر بوخت و، چې د علمي فکر لرونکو
فيلسوفانو يوه ډله وه او کارناپ هم پکې شامل
و۔ د کارناپ کار په ځانګړي ډول د ګوډل د
نتيجو تر اغېز لاندې و۔

په 1938 کښې ګوډل له اډېلې نيمبورسکي سره واده
وکړ۔ مور او پلار ئې خوشحاله نۀ وو: نۀ يوازې
دا چې هغه تر ګوډل شپږ کاله مشره وه او له
وړاندې ئې طلاق اخيستے و، بلکې په يوه شپني
کلب کښې ئې د نڅاګرې په توګه کار کاوۀ۔ خو
ټولنيزو فشارونو ګوډل اغېزمن نۀ کړ، او تر
وفاته ئې خوشحاله واده‌وال پاتې شول۔

وروسته له دې چې نازي جرمني په 1938 کښې
اتريش له ځان سره يوځاے کړ، ګوډل او اډېلې
متحده ايالاتو ته کډوال شول۔ هلته ئې د
نيو جرسي په پرنسټن کښې د لوړو څېړنو په
انسټيټيوټ کښې دنده واخيسته۔ له خپل درونګرا
او غيرمعمول طبيعت سره سره په پرنسټن کښې
د ګوډل وخت له ګډ کاره ډک او ګټور و۔ هغه
د سټونو په تيورۍ، فلسفه او فزيک کښې مقالې
خپرې کړې۔ په ځانګړي ډول ئې په IAS کښې
له خپل همکار البرټ اينشټاين سره ډېره
پياوړې ملګرتيا جوړه کړه۔

په وروستيو کلونو کښې د ګوډل ذهني روغتيا
خرابه شوه۔ په 1977 کښې د مېرمنې روغتون ته
تلل دا معنا لرله چې نور ئې هغه دپاره
خواړۀ نۀ شول پخولے۔ هغه چې د ټول ژوند
په اوږدو کښې له ذهني روغتيايي ستونزو سره
مخ و، په وهمي بدګمانۍ اخته شو۔ د زهر ورکولو
له سختې وېرې ګوډل له خوراکه انکار وکړ۔
هغه د لوږې له امله د 1978 د جنورۍ پر
14مه په پرنسټن کښې وفات شو۔

\begin{reading}
د ګوډل د ژوند د بشپړ ژوندليک دپاره
\citet{Dawson1997} وګورئ۔ د نورو ژوندليکي
ليکنو، او منطق او فلسفې ته د ګوډل د
مرستو په هکله د مقالو دپاره
\citet{Wang1990}، \citet{Baaz2011}،
\citet{Takeuti2003} او \citet{Sigmund2007}
وګورئ۔

د ګوډل د دوکتورا رساله په اصلي جرمني ژبه
شته \citep{Godel1929}۔ د ناتکميلۍ د قضيو
اصلي متن \citep{Godel1931} دے۔ د ګوډل
ټولې خپرې شوې او ناخپرې ليکنې، او د
ليکونو يوه ټاکلې برخه، په انګرېزي ژبه
د هغه په \emph{راټولو شويو ليکنو}
کښې شته \citet{Godel1986,Godel1990}۔

د ګوډل د ناتکميلۍ د قضيو د مفصل بيان
دپاره \citet{Smith2013} وګورئ۔ د ګوډل
د قضيو د غيرصوري، فلسفي بحث دپاره
د مارک لينسېنماير پوډکاسټ وګورئ
\citep{Linsenmayer2014}۔
\end{reading}

\end{document}
